\chapter{Introduction}
\source{cheeger-gromov-taylor:page-0001,cheeger-gromov-taylor:page-0002,cheeger-gromov-taylor:page-0003,cheeger-gromov-taylor:page-0004}

Let $M^n$ be a complete Riemannian manifold. The operator
$\Delta=-\delta d$ on functions is nonpositive and essentially self-adjoint on
$C_c^\infty(M)$. For suitable $f$, the spectral theorem defines
\begin{equation}
f(\sqrt{-\Delta})=\int_0^\infty f(\lambda)\,dE_\lambda,
\tag{0.1}
\end{equation}
where $dE_\lambda$ is the projection-valued measure of $\sqrt{-\Delta}$.
The operator has a kernel $k_{f(\lambda)}(x_1,x_2)$. In particular, the heat
kernel is $E(x_1,x_2,t)=k_{e^{-\lambda^2t}}$; for $t>0$ it is positive, symmetric,
and smooth in $(x_1,x_2,t)$, and for fixed $(x_2,t)$ belongs as a function of
$x_1$ to the domain of every positive power of $\Delta$.

For a domain $D\subset M^n$, write $\lVert\cdot\rVert_\infty$ for the supremum
norm and $\lVert\cdot\rVert_{L^2(D)}$ for the $L^2$ norm. The Sobolev inequality
has an optimal constant $S_D>0$ satisfying
\begin{equation}
S_D\left(\int_D|g|^{n/(n-1)}\right)^{(n-1)/n}
\leq \int_D\lVert dg\rVert,
\tag{0.2}
\end{equation}
for every $g\in H_0^1(D)$. If $u$ solves
$(\Delta-\partial_t)u=0$ on $[0,T]\times B_R(x)$, then
\begin{equation}
|u(x,T)|\leq c_n\left[T^{-(n+2)/2}+A^{-(n+2)}\right]^{1/2}S_D^{-n/2}
\lVert u\rVert_{L^2([0,T]\times B_A(x))}.
\tag{0.3}
\end{equation}
If $u$ satisfies $(\partial_t^2+\Delta)u=0$ on
$[-A/2,A/2]\times B_A(x)$, then
\begin{equation}
|u(x,0)|\leq c_nA^{-(n+1)/2}S_D^{-n/2}
\lVert u\rVert_{L^2([-A/2,A/2]\times B_A(x))}.
\tag{0.3'}
\end{equation}

The optimal Sobolev constant equals the optimal isoperimetric constant
$\Phi(D)$, characterized by
\begin{equation}
A(\partial U)\geq \Phi(D)V(U)^{(n-1)/n}
\tag{0.4}
\end{equation}
for every smooth subdomain $U\subset D$. Define $\widetilde\omega(D)$ by
letting $x\in D$, $v\in T_xM^n$ with $\lVert v\rVert=1$, letting $\gamma_v$ be
the geodesic with $\gamma_v'(0)=v$, and letting $t_v$ be the first parameter
for which $\gamma_v(t_v)\in\partial D$:
\begin{equation}
\widetilde\omega(D)=\inf_{x\in D}\mu\{v\in T_xM^n:\lVert v\rVert=1,
\gamma_v|_{[0,t_v]}\text{ is minimal}\}.
\tag{0.5}
\end{equation}
Here $\mu$ is normalized angular measure on the unit sphere. Croke's estimate is
\begin{equation}
\Phi(D)\geq 2^{(n-1)/n}\frac{\alpha(n-1)}{\alpha(n)^{(n-1)/n}}
\bigl(\widetilde\omega(D)\bigr)^{(n+1)/n},
\tag{0.6}
\end{equation}
where $\alpha(n)$ is the volume of the unit $n$-sphere; the estimate is sharp.

The heat-kernel estimates reduce to controlling an $L^2$ norm and a lower bound
for $\widetilde\omega(B_{r_2}(x_2))$. Finite propagation speed for
$\cos\sqrt{-\Delta}$ gives the required $L^2$ estimates and, via Fourier synthesis,
extends them to kernels $k_f$ with $f\in\mathcal F_2^{\varphi,A}$.

Define the injectivity angle by
\begin{equation}
\widetilde\omega(r,x)=\mu\{\gamma'(0):\gamma(0)=x,\ \gamma|_{[0,r]}
\text{ is minimal}\}.
\tag{0.7}
\end{equation}
Then
\begin{equation}
\widetilde\omega(B_r(x_0))\geq\inf_{x\in B_r(x_0)}\widetilde\omega(2r,x),
\tag{0.8}
\end{equation}
and, if $r\leq i(x)/2$ for all $x\in B_r(x_0)$,
\begin{equation}
\widetilde\omega(B_r(x_0))=\inf_{x\in B_r(x_0)}\widetilde\omega(2r,x)=1.
\tag{0.9}
\end{equation}
A lower sectional-curvature bound, an upper sectional-curvature bound, and a
lower bound for $i(p)$ yield a lower bound for $i(x)$ depending on
$\rho(p,x)$; even for negative lower curvature, this bound decays at most
exponentially. Alternatively, a lower Ricci bound and
$V(B_{r_0}(p))=V_{r_0}(p)\geq V>0$ directly give a lower bound for
$\widetilde\omega(2r,x)$. Consequently, assuming
$\Ric_M\geq(n-1)H$ and $V_r(p)\geq V>0$ gives upper estimates for the heat
kernel and for $k_f$ with $f\in\mathcal F_2^{\varphi,A}$.

If $i(x)$ has a lower bound and $|K_M|$ is bounded, elliptic estimates for a
parametrix of $\Delta$ give kernel estimates for the wider class
$f\in\mathcal F_1^N$, $N\geq n/4$. Under only
$V(B_{r_0}(p))\geq V>0$, one also obtains an estimate for the decay of $i(x)$.
In dimensions $2$ and $3$, the same class $\mathcal F_1^N$ is controlled under
only $\Ric_M\geq(n-1)H$ and $V_{r_0}(p)\geq V>0$.

If $\Ric_x\geq c/r^2$ outside $B_{r_0}(p)$, the volume $V_r(p)$ has sharp
asymptotic upper and lower bounds depending on $c$; its qualitative growth
changes with the value of $c$.
